\documentclass[11pt,a4paper]{article}
\usepackage[margin=1in]{geometry}
\usepackage{amsmath,amssymb}
\usepackage{parskip}
\usepackage{hyperref}

\title{\textbf{Pre-Experiment Overview:\\ Fine Structure of Two-Electron Atom He}}
\author{IBRAHIM H.I. ABUSHAWISH \\ Istanbul University, Department of Physics}
\date{}

\begin{document}
\maketitle

\section*{Purpose of the Experiment}
According to the lab manual, the main objective of this lab is to investigate the optical emission spectrum of the two-electron helium atom (He-I). By dispersing the emitted light through a high-resolution optical spectrometer equipped with a diffraction grating, we measure the diffraction angles for each visible spectral line, determine their precise wavelengths, identify singlet and triplet state transitions, and construct the spectrometer calibration curve.

\section*{Importance and Significance}
Studying helium is a crucial bridge in quantum mechanics: while hydrogen is a solvable one-electron system, helium introduces multi-electron electrostatic repulsion and quantum mechanical electron exchange symmetry. Because electrons are indistinguishable fermions, the total wavefunction must be anti-symmetric under particle exchange, giving rise to two distinct term systems: parahelium with antiparallel spins (singlet states, total spin $S = 0$) and orthohelium with parallel spins (triplet states, total spin $S = 1$). Because dipole selection rules strictly forbid transitions between different spin multiplicities ($\Delta S = 0$), helium essentially behaves optically as two separate, non-combining systems.

In the lab setup, a helium discharge tube emits bright spectral lines spanning red, yellow, green, blue, and purple. A spectrometer consisting of a collimator, prism/grating table, and telescope allows us to measure left and right diffraction angles $\theta_{\text{left}}$ and $\theta_{\text{right}}$ for orders $m = 1$ and $m = 2$. Using the grating relation $m \lambda = d \sin\theta_{\text{ave}}$, we calculate the experimental wavelengths of individual transitions (such as $706.5\text{ nm}$, $587.6\text{ nm}$, $501.5\text{ nm}$, $492.2\text{ nm}$, $471.3\text{ nm}$, and $447.1\text{ nm}$). Plotting $\sin\theta_{\text{ave}}$ against wavelength produces an accurate linear calibration curve and allows us to verify Russell-Saunders ($L$-$S$) coupling and the fine-structure term diagram of multi-electron atoms.

\end{document}
